\section{Cálculo rigoroso e reprodutibilidade}\label{sec:computation}

Esta seção descreve em que consistem as partes da prova assistidas por computador, em que é preciso confiar e como
um leitor pode conferi-las.

\subsection{Aritmética e base de confiança}\label{sec:computation-arith}

Três tipos de cálculo entram nos certificados.
\begin{enumerate}[label=(\alph*),leftmargin=*]
\item \emph{Aritmética racional exata} (o módulo \texttt{fractions} do Python): as desigualdades finais de todo
  certificado, a saber, as condições de Perron $Nv\le\theta v$ com $\theta<1$ do
  \cref{lem:perron}, as condições de Newton--Kantorovich da \cref{sec:roots-nk}, as coberturas dos contornos e
  da região da \cref{prop:Kinj}, o polinômio de deflação $q$ e o fato F1.
\item \emph{Aritmética de bolas} com Arb~\cite{Arb}, por meio do \texttt{python-flint}: as cotas dos símbolos do
  \cref{lem:F3} e a folga entre os pontos da grade.
\item \emph{Cotas a posteriori sobre cálculos em ponto flutuante.} A álgebra linear de grande porte (dimensão $14016$ para a
  referência finita, janelas de dimensão até $11200$) é feita em dupla precisão IEEE. Os resultados em ponto flutuante
  são usados apenas como inversas aproximadas, resolvedores ou vetores de teste, e toda quantidade que entra num certificado é
  cotada rigorosamente: os resíduos são calculados com cotas de arredondamento entrada a entrada da forma padrão
  $\gamma_n\||M||X|\|$~\cite{Higham2002}; as normas de inversas são cotadas por desigualdades de Loewner verificadas,
  $t^2MM^*-I\succeq0$, cuja positividade definida é verificada pelo método de Rump~\cite{Rump2006} aplicado à
  realificação; e os fatores de acoplamento são certificados como na \cref{sec:counting-reference}.
\end{enumerate}
A base de confiança é, portanto: a correção das cotas publicadas em~\cite{RT2019}, que usamos tal como enunciadas;
a aritmética IEEE~754 com arredondamento para o mais próximo, como supõem as cotas a posteriori; o Arb; e o código Python do
programa de verificação, que é curto para cada conferência e está disponível para inspeção. Nenhuma parte da cadeia final
está formalizada num assistente de provas.

Como as entradas em ponto flutuante dependem da biblioteca BLAS e do número de threads, refazer um
certificado em outra máquina reproduz o $\theta$ racional apenas a menos de $\sim10^{-16}$, não exatamente; o
critério de sucesso de uma reprodução é $\theta<1$, com uma diferença na escala do arredondamento.

\subsection{Os certificados}\label{sec:computation-certs}

A \cref{tab:certs} lista os certificados calculados, os resultados que eles estabelecem e o seu custo aproximado. Os
cálculos foram feitos em cinco computadores comuns, com quatro a seis núcleos e $7$--$31$\,GB de memória; a cadeia
inteira leva alguns dias em três ou quatro dessas máquinas. Os maiores cálculos individuais são a decomposição de Schur
da referência finita e as cotas de Loewner do seu resolvente (horas por centro do contorno), e as
caudas com $26$ janelas densas (três a seis horas por centro).

\begin{table}[ht]
\centering
\small
\begin{tabular}{@{}p{0.29\textwidth}p{0.35\textwidth}p{0.27\textwidth}@{}}
\toprule
certificado & estabelece & custo \\
\midrule
símbolos de $B$ e $B^\sharp$ (Arb, grades $2048\times4096$ e $1024\times2048$) & \cref{lem:F3}; \cref{thm:counts}~(i) & horas, 4 processos \\
recuperação e transferência do raio & \cref{lem:Lreal,lem:recovery} & minutos \\
referência finita: matriz, Schur, resíduo & \cref{lem:schur} & horas \\
nível $Z$ em $174$ discos (Loewner, resoluções de Schur) & \cref{thm:counts}~(ii),(iii) & $\sim$9 min por ponto \\
fatores de acoplamento de posto baixo ($13$ centros) & \cref{sec:counting-reference} & $\sim$1 h por centro \\
caudas, $26$ janelas ($13$ centros) & \cref{sec:counting-reference} & 3--6 h por centro \\
janelas dentro das faixas & \cref{lem:windows} & horas \\
NK no sistema aumentado em $\mu,\sK,\sphys,\sB$ & \cref{prop:nk,prop:mu} & 3--6 h cada \\
testemunhos dos vínculos & \cref{prop:witness} & minutos \\
$154$ células e Rouch\'e para $K$ & \cref{prop:Kinj} & $\sim$1 dia, 3 máquinas \\
F1 & \cref{lem:F1} & segundos \\
\bottomrule
\end{tabular}
\caption{Certificados calculados. Os nomes dos arquivos e os hashes sha256 estão listados no \cref{app:certificates}.}
\label{tab:certs}
\end{table}

\subsection{O programa de verificação}\label{sec:computation-verify}

O repositório contém um programa de verificação com dois níveis.

O \emph{nível 1} refaz, em aritmética racional exata e com o código do repositório, as desigualdades finais
dos certificados a partir dos arquivos de certificado gravados (pequenos arquivos JSON) e dos dados publicados de~\cite{RT2019}:
\begin{itemize}[leftmargin=*]
\item a cota de Perron em três níveis em cada um dos $175$ pontos dos $22$ arquivos de discos das duas faixas ($174$ discos de
  raio positivo, e uma entrada de raio zero que não é usada na cobertura);
\item as coberturas de $\partial\Omega_A$ e $\partial\Omega_B$;
\item as contagens finitas;
\item as cotas das janelas;
\item os quatro certificados de Newton--Kantorovich, os testemunhos e a identificação da raiz de gauge;
\item as $154$ células e a sua cobertura da região da \cref{prop:Kinj} fora do disco $|s-0.401|<0.05$ do
  argumento de Rouch\'e para $K$;
\item o polinômio de deflação $q$ da \cref{sec:counting-rouche};
\item as cotas de recuperação e de transferência do raio, a cota $\beta_+$ da parte do fundo e o F1.
\end{itemize}
Alguns certificados entram no nível 1 apenas pelos seus valores gravados, e só são recalculados no nível 2: as cotas dos
símbolos e das partes livres (o nível 1 recombina as faixas gravadas), e as cotas de Perron das $16$ células do
argumento de Rouch\'e para $K$.
Ele leva cerca de duas horas em quatro núcleos. Todas as conferências passam: a pior raiz de Perron nos discos é $0.999074$, com
diferenças para os valores gravados $\le7.8\cdot10^{-16}$, e, nas máquinas em que foram gerados, os dados de
Newton--Kantorovich, os testemunhos e as células são reproduzidos de forma idêntica. Em outras plataformas, as
entradas em ponto flutuante dos certificados podem diferir nos últimos bits, e o mesmo acontece com os racionais recalculados; o programa
confere então que cada certificado recalculado é válido, não passa das cotas publicadas aqui e concorda com o gravado
a menos de $10^{-9}$ relativo.

O \emph{nível 2} recalcula os certificados intermediários a partir dos dados de~\cite{RT2019}. Como amostra, recalculamos
do zero, numa máquina diferente:
\begin{itemize}[leftmargin=*]
\item a matriz de referência, idêntica bit a bit à gravada;
\item a cota do resíduo de Schur, idêntica;
\item a cauda do centro de contorno $3.0$ ($26$ janelas e o nível distante, $5.3$ horas), idêntica em todos os campos.
\end{itemize}

\subsection{Disponibilidade de dados e código}\label{sec:computation-data}

O código, os arquivos de certificado, o programa de verificação e um guia de reprodução estão disponíveis em
\url{https://github.com/edivaldo-amaral/choptuik-espectro} (licença MIT)~\cite{Repo2026}. Os arquivos de dados grandes, a saber, a matriz
de referência, os seus fatores de Schur e os fatores de acoplamento ($16$\,GiB), estão depositados no Zenodo, DOI
\href{https://doi.org/10.5281/zenodo.23222363}{10.5281/zenodo.23222363} (CC BY 4.0)~\cite{Zenodo2026}, com hashes sha256 que
também estão registrados nos arquivos de certificado. Os dados de~\cite{RT2019} não são redistribuídos; um script os baixa
do fonte de~\cite{RT2019} no arXiv e confere os seus hashes sha256.

\subsection{Conferências independentes assistidas por IA}\label{sec:computation-reviews}

Antes de serem usados, os certificados e os lemas foram conferidos por agentes de IA independentes dos que produziram o
trabalho, com as exceções listadas abaixo. Essas conferências não são revisões por especialistas humanos. O protocolo foi
desenhado para que uma conferência produza evidência verificável, e não uma opinião:
\begin{itemize}[leftmargin=*]
\item o agente conferente rodava num modelo diferente do que produziu o trabalho, com contexto limpo;
\item ele recebia as alegações enunciadas com precisão, as derivações sob revisão, o código e os dados, mas não
  as conclusões do autor nem auditorias anteriores;
\item ele tinha de entregar um artefato por alegação: uma derivação independente, uma implementação independente ou um
  cálculo direto da quantidade verdadeira para comparar com a cota;
\item antes de cada conferência era tirado um snapshot com manifesto sha256, comparado ao final dela.
\end{itemize}
Houve doze conferências desse tipo:
\begin{itemize}[leftmargin=*]
\item as células da \cref{prop:Kinj} (uma conferência anterior delas, com outra ferramenta, foi superficial e não é
  contada);
\item os certificados em $\sK$ (com o argumento de Rouch\'e para $K$), em $\mu$ e em $\sphys$;
\item os argumentos de Rouch\'e das duas faixas, incluindo o certificado em $\sB$;
\item os lemas estruturais e a cadeia da prova;
\item o lema de gauge, a reconstrução direta e a recíproca;
\item as hipóteses de modelagem;
\item o manuscrito inteiro, alegação por alegação, contra os arquivos de certificado.
\end{itemize}
Uma parte não foi conferida isoladamente: os cálculos dos símbolos do \cref{lem:F3} foram examinados apenas no seu uso,
não recalculados. Nenhuma das conferências derrubou um resultado. Todas encontraram lacunas, sobretudo de rastreabilidade ou de enunciado,
e algumas de derivação: por exemplo, uma encontrou que uma versão inicial das células usava a aproximação diádica de
$\mu$ em vez do valor verdadeiro, e outra que o erro da deflação fora da caixa finita tinha sido omitido. Todas
foram corrigidas antes de os resultados serem usados; os relatórios e artefatos estão no repositório.

\subsection{Uso de assistentes de IA}\label{sec:computation-ai}

Este trabalho foi feito com uso extensivo de assistentes de IA, entre setembro e outubro de 2026. Descrevemos o seu papel
como exige a política da revista.

\emph{Escopo.} Os assistentes escreveram a maior parte do código, fizeram a maior parte das derivações e dos cálculos, prepararam
a documentação e redigiram este manuscrito. A sua contribuição é substancial e não se limita à edição do
texto.

\emph{O papel do autor.} O autor propôs o problema, o espectro instável da solução de~\cite{RT2019}, que
os seus autores deixaram em aberto; dirigiu o projeto e escolheu entre as abordagens; rodou os cálculos nas suas próprias
máquinas e em máquinas emprestadas; conferiu os argumentos e os resultados; e assume total responsabilidade pelo conteúdo
deste artigo.

\emph{Ferramentas.} O assistente principal foi o Claude (Anthropic), usado por meio do ambiente Claude Code, com os modelos
\texttt{claude-opus-5} e, depois, \texttt{claude-opus-5-5}. O agente Codex, da OpenAI, foi usado na primeira sessão
exploratória e durante um intervalo de três dias, com os modelos \texttt{gpt-5.6-luna}, \texttt{gpt-5.6-sol},
\texttt{gpt-6-astra} e \texttt{gpt-reserve}; tudo o que foi produzido nesse intervalo foi auditado e rederivado ou
conferido antes do uso. As conferências da \cref{sec:computation-reviews} foram feitas por agentes Claude rodando
\texttt{claude-fable-5-1}, um modelo diferente do que produziu o trabalho. A pedido do autor, uma versão preliminar do
manuscrito também foi lida pelo ChatGPT (OpenAI, modelo Sol 6.1); os seus comentários foram conferidos contra os certificados e
as fontes antes de serem incorporados. Os identificadores dos modelos são os registrados nos logs das sessões; um relato datado das sessões está no arquivo \texttt{paper/AI\_USE.md} do
repositório. Esta tradução para o português também foi feita com o Claude (\texttt{claude-opus-5-5}) e revisada pelo autor.

\emph{Limites da verificação.} Os argumentos analíticos estão neste artigo; algumas estimativas de rotina aparecem em
forma de esboço, com as suas constantes conferidas por código no repositório (\cref{app:lemmas}). Toda cota calculada é
produzida por código que pode ser inspecionado e rodado de novo a partir dos arquivos de certificado e dos dados. As desigualdades finais dos
certificados são refeitas pelo programa de verificação (nível~1 da \cref{sec:computation-verify}); algumas cotas
intermediárias entram no nível~1 apenas pelos seus valores gravados, como listado na \cref{sec:computation-verify}, e só são recalculadas
no nível~2, no qual uma amostra dos certificados intermediários foi recalculada a partir dos dados
de~\cite{RT2019}. Os lemas e a maior parte dos certificados também foram
examinados pelas conferências da \cref{sec:computation-reviews}, com as exceções listadas ali. Nada disso substitui a
revisão por pares; o objetivo é que a confiança no resultado não dependa do processo que o produziu.
